\begin{abstract}
% 向量量化通过联合编码多个权重提高低比特表示能力，而高维结构化方法通常通过扩大有效编码维度进一步提升量化质量。本文提出一个互补视角：基础码字维度与码字选择的联合范围是两个独立的设计维度。在共享尺度 VQ 中，尺度与多个码字选择相互耦合，使逐点更新可能陷入停滞，而直接联合搜索又具有指数级复杂度。为此，我们提出共享尺度协同码字优化方法，通过多尺度候选生成、残差感知的可分离二次上界和直线下包络搜索，将指数级组合搜索转化为有限代表区间上的硬索引与尺度联合求解。实验表明，联合优化将 Qwen3-4B/8B 的 WikiText-2 PPL 从 22.69/14.41 降至 12.18/9.76，显著优于仅优化尺度；在 Qwen3.8-27B 上，本方法相较 QTIP 和 LLVQ 获得更低的 WikiText-2/C4 PPL，并在 LiveCodeBench v6 和 GPQA-D 上取得更高分数。结果表明，扩大联合决策范围能够在保持低维码字表示的同时显著提升低比特 VQ 的量化质量。
Vector quantization improves low-bit representations by jointly encoding multiple weights, while high-dimensional structured methods typically improve quantization quality by increasing the effective coding dimension. We offer a complementary perspective: the base-codeword dimension and the joint range of codeword selection are independent design axes. In shared-scale VQ, the scale is coupled with multiple codeword choices, so pointwise updates may stagnate, whereas direct joint search has exponential complexity. We therefore propose shared-scale cooperative codeword optimization, which combines multi-scale candidate generation, a residual-aware separable quadratic upper bound, and lower-envelope search to convert exponential combinatorial search into joint hard-index and scale optimization over finitely many representative intervals. Joint optimization reduces WikiText-2 perplexity on Qwen3-4B/8B from 22.69/14.41 to 12.18/9.76, substantially outperforming scale-only optimization. On Qwen3.8-27B, our method obtains lower WikiText-2/C4 perplexity than QTIP and LLVQ and higher scores on LiveCodeBench v6 and GPQA-D. These results show that expanding the joint decision range can substantially improve low-bit VQ while retaining low-dimensional codewords.
\end{abstract}

% 关键词：向量量化；共享尺度；码字联合选择；二次上界；下包络搜索。
\noindent\textbf{Keywords:} vector quantization; shared scale; joint codeword selection; quadratic upper bound; lower-envelope search
